\begin{table}[H]
  \caption{Trajectory-level local starting-coordinate sensitivity at four
  sites identified by the structural-concordance analysis.}
  \label{tab:local-coordinate-sensitivity}
  \centering
  \scriptsize
  \setlength{\tabcolsep}{3pt}
  \renewcommand{\tabularxcolumn}[1]{m{#1}}
  \begin{tabularx}{\linewidth}{@{}>{\raggedright\arraybackslash}m{2.2cm}>{\raggedright\arraybackslash}m{2.5cm}S[table-format=1.3]cS[table-format=1.3]S[table-format=1.3]S[table-format=+1.3,retain-explicit-plus=true]>{\raggedright\arraybackslash}X@{}}
    \toprule
    Target & Metric & \multicolumn{1}{c}{\shortstack{Initial model--\\reference RMSD\\(\si{\angstromunit})}} &
    \shortstack{Reference\\overlap ($n/21$)} & \shortstack{Median overlap\\fraction} &
    \multicolumn{1}{c}{\shortstack{Median experimental-\\closer fraction}} &
    \multicolumn{1}{c}{\shortstack{$\Delta$ experimental-\\reference RMSD\\(\si{\angstromunit})}} &
    \shortstack[l]{Event-fraction contrast\\(95\% interval)} \\
    \midrule
    TfCut1 H185 & Side-chain conformation & 1.509 & 21/21 & 1.000 & 1.000 & -1.036 & Not estimable ($n=0$) \\
    FoCut5a L184 & Local backbone & 1.543 & 1/21 & 0.010 & 0.041 & +0.813 & \shortstack[l]{Heavy-atom contact ($n=11$):\\$-0.071$ ($-0.150$, $-0.001$)} \\
    HiC T164 & Side-chain conformation & 2.076 & 16/21 & 0.077 & 0.147 & -0.200 & \shortstack[l]{Heavy-atom contact ($n=13$):\\$+0.031$ ($-0.035$, $+0.096$)} \\
    HiC T166 & Side-chain conformation & 2.053 & 20/21 & 0.553 & 0.820 & -0.825 & \shortstack[l]{Side-chain H bond ($n=13$):\\$-0.008$ ($-0.072$, $+0.053$)} \\
    \bottomrule
  \end{tabularx}

  \vspace{4pt}
  \parbox{\linewidth}{\footnotesize
  The initial model--reference RMSD is the RMSD between the production-model
  starting conformation and the experimental-reference conformation. Reference
  overlap is the number of trajectories in which
  at least 5\% of the 20--100~ns frames met the corresponding diagnostic
  threshold. Median overlap fraction is the median trajectory-level overlap
  fraction. Median experimental-closer fraction is the median fraction of
  analysis frames with a lower RMSD to the experimental-reference state than
  to the production-model state.
  For each trajectory, $\Delta$ experimental-reference RMSD was calculated as its
  mean RMSD to the experimental reference over 20--100~ns minus its RMSD to that
  reference at 0~ns. The table reports the unweighted arithmetic mean of these
  differences across all 21 trajectories for each target. Negative values
  indicate closer agreement during the analysis window.
  In the last column, $n$ denotes the number of retained trajectories containing
  frames in both closer-state classes. Contrasts are the arithmetic means of
  the within-trajectory experimental-closer minus production-closer event-fraction
  differences. Percentile 95\% intervals were obtained from 20,000 bootstrap
  resamples of these $n$ trajectories.
  For H185, each retained trajectory contained frames from only one closer-state
  class, so a within-trajectory contrast could not be estimated.
  Values are rounded to three decimal places.}
\end{table}
